\section{Day 4: Number Fields (Sep.\ 17, 2026)}
Office hours have been moved to 2--3pm on Thursdays. Last time, given an algebraic number $\alpha \in \QQalg$, we defined the conjugates of $\alpha$ to be the algebraic numbers with the same minimal polynomial, and we introduced the following lemma,
\begin{lemma}
    \begin{parlist}
        \item Let $f \in \QQ[x]$ be monic and irreducible of degree $d$; then $f$ has $d$ distinct roots in $\CC$.
        \item If $\alpha$ is an algebraic number of degree $d$, then there exist exactly $d$ algebraic numbers conjugate to $\alpha$.
    \end{parlist}
\end{lemma}
\begin{proof}
    It is clear that (i) implies (ii), so it suffices to prove (i). First, $\CC$ is algebraically closed, so $f = (x - \alpha_1) \dots (x - \alpha_d)$ with each $\alpha_i \in \CC$. We need to show that $\alpha_i \neq \alpha_j$ for all $i, j$; suppose $\alpha \in \CC$ is a root of $f$ with multiplicity greater than $1$. Then $f$ being monic and irreducible implies that $f$ is the minimal polynomial of $\alpha$. By taking the derivative of $f$, we see that $f'(\alpha) = 0$; but then $\deg f' = d - 1$, which is a contradiction.
\end{proof}

Let's do some examples.
\begin{enumerate}[(i)]
    \item The conjugates of $\sqrt{5}$ are $\sqrt{5}$ and $\sqrt{-5}$, since $x^2 - 5 = 0$ is the minimal polynomial of $\sqrt{5}$. Recall that it is easy to check if a quadratic polynomial is irreducible (check the rational roots).
    \item Let $f(x) = x^4 - 10x^2 + 1$, which we know (from previous classes) may be factored as
    \[ f(x) = (x - \sqrt{2} - \sqrt{3})(x - \sqrt{2} + \sqrt{3})(x + \sqrt{2} - \sqrt{3})(x + \sqrt{2} + \sqrt{3}). \]
    % We saw that, if $\alpha$ satisfies $f(\alpha) = \alpha^4 - 10 \alpha^2 + 1 = 0$, then 
    We claim that $f$ is irreducible; there are many ways to check this, see the MAT347 notes. {\color{airforceblue} In my opinion, the easiest would be to check the extensions of $\QQ$.}
    \begin{corollary}[Vi\`ete] \label{cor:viete_formulas}
        If $\alpha \in \QQalg$ has minimal polynomial $f(x) = x^d + a_{d-1} x^{d-1} + \dots + a_0$ with conjugates $\alpha_1, \dots, \alpha_d$, then 
        \[ \alpha_1 \dots \alpha_d = (-1)^d a_0, \quad \alpha_1 + \dots + \alpha_d = -a_{d-1}. \]
    \end{corollary}
    The proof is immediate; just write $f(x) = (x - \alpha_1) \dots (x - \alpha_d)$.
\end{enumerate}

\begin{definition} \label{defn:number_field}
    An \textit{algebraic number field} is a subfield of $F$ of $\CC$ which is finite-dimensional over $\QQ$.
\end{definition}
Equivalently, an algebraic number field is a finite extension of $\QQ$. (recall from MAT347 that any finite extension is generated by a single element, i.e., the primitive element theorem. cf.\ D\&F \S14.4.25).
If $F$ is a number field, then the \textit{ring of integers} of $F$ is denoted $\SO_F = F \cap \ZZalg$.

\begin{remark} \label{rmk:properties_of_number_fields}
    We have the following properties of number fields.
    \begin{enumerate}[(i)]
        \item Let $V = F \subset \CC$ be a number field; then $\alpha \cdot V \subset V$ for $\alpha \in F$, and so $\alpha \in \QQalg$. Specifically, any element of an algebraic number field is an algebraic number, $F \subset \QQalg$.
        \item $\SO_F$ is a subring of $\CC$ containing $\ZZ$.
        \item $\QQ \cap \ZZalg = \ZZ$, i.e., $\SO_\QQ = \ZZ$.
        \item Let $\alpha \in \QQalg$; we have that $\QQ(\alpha) \subset \CC$ is an algebraic number field. The primitive element theorem states that \textit{every} number field is of this form.
    \end{enumerate}
\end{remark}

\newpage
\begin{example}
    Let $d \in \ZZ \setminus \{0, 1\}$ be squarefree. Let $F = \QQ(\sqrt{d}) = \{a + b \sqrt{d} \mid a, b \in \QQ\}$. What is $\SO_F$? Observe that
    \[ (a + b \sqrt{d})(a - b \sqrt{d}) = a^2 - db^2, \quad (a + b \sqrt{d}) + (a - b \sqrt{d}) = 2a, \]
    so the minimal polynomial is $x^2 - 2ax + (a^2 - db^2)$. We want that $2a \in \ZZ$ and $a^2 - db^2 \in \ZZ$. If $d \equiv 2, 3 \pmod{4}$, then $a, b$ are actually integers (checked by bashing). If $d \equiv 1 \pmod{4}$, then $a, b$ are both integers, or half-integers. Thus,
    \[ \SO_F = \begin{cases} \ZZ + \ZZ \sqrt{d}, &\quad d = 2, 3 \pmod{4}, \\
        \ZZ + \ZZ\!\left(\!\frac{-1 + \sqrt{d}}{2}\!\right),& \quad d = 1 \pmod{4}. \end{cases} \]
\end{example}
Note that, for $\alpha \in \ZZalg$ and $F = \QQ(\alpha)$, $\SO_F$ can be bigger than $\ZZ[\alpha]$.

\begin{lemma} \label{lem:field_of_fractions}
    \begin{enumerate}[(i)]
        \item If $\alpha \in \QQalg$, there exists a nonzero integer $m \in \NN$ such that $m\alpha \in \ZZalg$.
        \item If $F$ is a number field, then $F = \Frac(\SO_F)$.
    \end{enumerate}
\end{lemma}
\begin{proof}
    We start by proving that (i) implies (ii). Since $\ZZ \subset \SO_F$, we have that $\alpha = \frac{m \alpha}{m} \in \Frac(\SO_F)$. {\color{airforceblue} Note that he skips a lot of intuition here, but the point is that $\SO_F$ contains the integers, so from (i), we may write $\alpha$ as a fraction in terms of elements of $\SO_F$; since $\SO_F = F \cap \ZZalg$, it is enough to have $\alpha \in \SO_F$ to get (ii).}

    To prove (i), let $f(x) = x^n + a_{n-1} x^{n-1} + \dots + a_0$ (with $a_i \in \QQ$) be the minimal polynomial of $\alpha$. Then there exists $m$ such that $m a_i \in \ZZ$ for all $i$; it follows that we may multiply $f(x)$ throughout by $m^n$ to see that $m \alpha$ is a root of
    \[ x^n + m a_{n-1} x^{n-1} + m^2 a_{n-2} x^{n-2} + \dots + m^n a_0 \in \ZZ[x]. \qedhere \]
\end{proof} \vspace{-8pt}

We do more examples now.
\begin{enumerate}[(i)]
    \item Consider the Gaussian integers $F = \QQ(i)$, where $i = \sqrt{-1}$ and $-1 \equiv 3 \pmod{4}$. Then $\SO_F = \ZZ[i]$, for which $\SO_F$ is a Euclidean domain with Euclidean norm $N(\alpha) = \alpha \ol \alpha$.
    \item Let $\omega = \exp(2\pi i / 3) \in \CC$. We have that $F = \QQ(\omega) = \QQ(\sqrt{-3})$. Since $-3 \equiv 1 \pmod{4}$, we have that $\SO_F = \ZZ + \ZZ \omega = \ZZ[\omega] \neq \ZZ[\sqrt{-3}]$. We claim that this is a Euclidean domain, with $N(\alpha) = \alpha \ol \alpha$ the Euclidean norm.
    \item Let $F = \QQ{\sqrt{-5}}$. $-5 \equiv 3 \pmod{4}$, so $\SO_F = \ZZ[\sqrt{-5}]$. Consider the Euclidean norm $N(a + b \sqrt{-5}) = a^2 + 5 b^2 \in \ZZ_{\geq 0}$. 
    \begin{lemma}
        The only units in $\ZZ[\sqrt{-5}]$ are $-1$ and $1$. The elements $3, 7, 1 + 2 \sqrt{-5}, 1 - 2 \sqrt{-5}$ are all irreducible.
    \end{lemma}
    \begin{proof}
        We claim that $\alpha \in \ZZ[\sqrt{-5}]$ is a unit if and only if $N(\alpha) = 1$. Suppose $\alpha \beta = 1$; then $N(\alpha \beta) = N(\alpha) N(\beta) = 1$. Since $N(\alpha) \in \ZZ_{\geq 0}$, if $N(\alpha) = 1$, then $\alpha \ol \alpha = 1$, meaning $\ol \alpha = \alpha\inv$. In this manner, it is enough to solve $a^2 + 5b^2 = 1$, from which we get $b = 0$ and $a = \pm 1$, whence the units are $\pm 1$.

        We now check that our claimed elements are irreducible. Suppose $3$ is not irreducible. Let us write $3 = \alpha \beta$, where $\alpha, \beta$ are non-units. Per above, we see that this must mean that their norms are not $1$; so
        \[ 9 = N(3) = N(\alpha) N(\beta) \implies N(\alpha) = N(\beta) = 3; \]
        since $\alpha = a + b \sqrt{-5}$ with $N(\alpha) = a^2 + 5b^2 \neq 3$, we get a contradiction, and we are done. The process is analogous to check that $7$ is irreducible, i.e., we arrive at a contradiction when we find that $a^2 + 5b^2 = 7$ has no integer solutions.

        \newpage
        To check that $1 \pm 2 \sqrt{-5}$ is irreducible (we'll only check the $1 + 2 \sqrt{-5}$ case because the other follows similarly), suppose $1 + 2 \sqrt{-5} = \alpha \beta$, where $\alpha, \beta$ are non-units. Since
        \[ 21 = N(1 + 2 \sqrt{-5}) = N(\alpha) N(\beta), \]
        we must have $N(\alpha)$ is $3$ or $7$, which is not possible as we saw earlier. Thus, we are done.
    \end{proof}
    {\color{airforceblue} The moral of this situation is that, to analyze Diophantine equations, we really just want to bash things out with factorization. In addition, norms are really useful for testing these properties as well!}
\end{enumerate}

We recall some important properties of field extensions (perhaps, from MAT347). Let $L/F$ be a finite extension (i.e., $F \subset L$). We denote the degree of the extension as
\[ \deg L/F = \abs{L : F} \coloneq \dim_F L. \]
As an example, $\abs{\QQ(\alpha) : \QQ} = \deg \alpha$. The tower law states that, if $L/M$ and $M/F$ are finite extensions, and $\omega_1, \dots, \omega_n$ is a basis for $L/M$, then
\[ L = M \omega_1 \oplus \dots \oplus M \omega_n. \]
We also have
\[ \abs{L : F} = \abs{L : M} \abs{M : F}. \]
We have the following key example. Let $F$ be a number field, and take $\alpha \in F$; then
\[ \abs{F : \QQ} = \abs{F : \QQ(\alpha)} \abs{\QQ(\alpha) : \QQ}, \]
for which the latter is simply just $\deg \alpha$. In particular, $\deg \alpha \mid \abs{F : \QQ}$. What can we do with this? Suppose $\alpha = \sqrt{2} + \sqrt{3}$. $\QQ(\alpha)$ contains $\sqrt{2}$ and $\sqrt{3}$. Since
\[ \sqrt{2} = \frac{\alpha^3 - 9 \alpha}{-4}, \]
we see that $\QQ(\sqrt{2}) \subset \QQ(\alpha)$. By writing
\[ \deg \alpha = \abs{F : \QQ} = \abs{F : \QQ(\sqrt{2})} \abs{\QQ(\sqrt{2}) : \QQ} = 4, \]
where $\lvert \QQ(\sqrt{2}) : \QQ \rvert = 2$, we see that $\lvert F : \QQ(\sqrt{2}) \rvert = 2$.

We now discuss norm and trace.
\begin{definition} \label{defn:norm_and_trace}
    Let $L/F$ be a finite extension. The \textit{norm} of $\alpha$, denoted $N_{L/F}(\alpha)$, is the determinant of the $F$-linear map $\phi_\alpha = \bullet \alpha \colon L \to L$. Similarly, the trace of $\alpha$ is $\Tr_{F/L}(\alpha) \coloneq \tr(\phi_\alpha)$.
\end{definition}

Let us look at some examples.
\begin{enumerate}[(i)]
    \item $\CC/\RR$ has basis $\{1, i\}$. The matrix for multiplication by $x + iy$ is
    \[ \begin{pmatrix} x & -y \\ y & x \end{pmatrix}, \]
    (which you may perhaps recognize as the Cauchy--Riemann equations). Then $N_{\CC/\RR}(x + iy) = x^2 + y^2$, and $\Tr_{\CC/\RR}(x + iy) = 2x$.

    \newpage
    \item $\QQ(i)/\QQ$ follows the exact same process. The norm and trace of $a + bi$ is $a^2 + b^2$, $2a$ respectively, since the basis and matrix are the same.
    
    \item $\QQ(\omega)/\QQ$, where $\omega^2 + \omega + 1 = 0$ (and $\omega = \exp(2\pi i / 3)$), has basis $1, \omega$. Multiplication by $a + b \omega$ in this basis is
    \[ \begin{pmatrix} a & -b \\ b & a - b \end{pmatrix}, \]
    so $N_{\QQ(\omega)/\QQ}(a + b\omega) = a^2 + b^2 - ab$, and $\Tr_{\QQ(\omega)/\QQ}(a + b \omega) = 2a - b$. The norm is the product of the conjugates, and conjuagte in this situation is the complex conjugate (really, this has to do with the fact that the degree is $2$; not every degree $2$ extension is like this, though).
\end{enumerate}

\begin{proposition} \label{prop:properties_of_norm}
    We give the main properties of the norm here. $N_{L/F} \colon L \to F$ has the following properties,
    \begin{enumerate}[(i)]
        \item $N_{L/F}(\alpha_1 \alpha_2) = N_{L/F}(\alpha_1) N_{L/F}(\alpha_2)$;
        \item $N_{L/F}(a) = a^{\abs{L:F}}$ if $a \in F$;
        \item $N_{L/F}(\alpha) = 0$ if and only if $\alpha = 0$;
        \item If $M$ is intermediate between $L$ and $F$, then
        \[ N_{L/F} = N_{M/F} \circ N_{L/M}. \]
    \end{enumerate}
\end{proposition}
\begin{proof}
    Properties (i) to (iii) are clear, so we will check (iv) only. We will start with a proof sketch in the special case that $\alpha \in M$; let $\omega_1, \dots, \omega_n$ be a basis for $L/M$. Then $L = M \omega_1 + \dots + M \omega_n$, and $\alpha$ acts block-diagonally. Thus, $N_{L/F}(\alpha) = N_{M/F}(\alpha)^n = N_{M/F} \circ N_{L/M}(\alpha)$. {\color{airforceblue} Perhaps, this will be continued later?}
\end{proof}

\begin{proposition} \label{prop:properties_of_trace}
    $\Tr_{L/F}$ has the following properties,
    \begin{enumerate}[(i)]
        \item $\Tr_{L/F}(\alpha_1 + \alpha_2) = \Tr_{L/F}(\alpha_1) + \Tr_{L/F}(\alpha_2)$;
        \item $\Tr_{L/F}(a \alpha) = a \Tr_{L/F}(\alpha)$ for $a \in F$ and $\alpha in L$;
        \item $\Tr_{L/F}(a) = \abs{L : F} a$ if $a \in F$;
        \item If $M$ is an intermediate field between $L$ and $F$, then
        \[ \Tr_{L/F} = \Tr_{M/F} \circ \Tr_{L/M}. \]
    \end{enumerate}
\end{proposition}
\begin{proof}
    The proof is the same.
\end{proof}

\begin{remark}
    Let $L/F$ be Galois. Then
    \[ N_{L/F}(\alpha) = \prod_{\sigma \in \Gal(L/F)} \sigma(\alpha), \quad \Tr_{L/F}(\alpha) = \sum_{\sigma \in \Gal(L/F)} \sigma(\alpha). \]
    In general, sum and products embed $L$ into $F$. {\color{airforceblue} I didn't catch what he said here; he'll probably restate this next class, so, we'll see then.}
\end{remark}
We will do this next time.